\documentclass[12pt,a4paper]{article}
\usepackage[utf8]{inputenc}
\usepackage{amsmath, amssymb, amsthm}
\usepackage{graphicx}
\usepackage{geometry}
\geometry{margin=1in}
\usepackage{hyperref}
\usepackage{listings}
\usepackage{xcolor}

\title{\textbf{Fourier Neural Operators (FNO) for Parametric PDEs}\\ \Large A Comprehensive Deep Dive}
\author{Battery Thermal Runaway Prediction Project}
\date{\today}

\definecolor{codegreen}{rgb}{0,0.6,0}
\definecolor{codegray}{rgb}{0.5,0.5,0.5}
\definecolor{codepurple}{rgb}{0.58,0,0.82}
\definecolor{backcolour}{rgb}{0.95,0.95,0.92}

\lstdefinestyle{mystyle}{
    backgroundcolor=\color{backcolour},   
    commentstyle=\color{codegreen},
    keywordstyle=\color{magenta},
    numberstyle=\tiny\color{codegray},
    stringstyle=\color{codepurple},
    basicstyle=\ttfamily\footnotesize,
    breakatwhitespace=false,         
    breaklines=true,                 
    captionpos=b,                    
    keepspaces=true,                 
    numbers=left,                    
    numbersep=5pt,                  
    showspaces=false,                
    showstringspaces=false,
    showtabs=false,                  
    tabsize=2
}
\lstset{style=mystyle}

\begin{document}

\maketitle
\tableofcontents
\newpage

\section{Theoretical Foundations and Context}

\subsection{The Challenge of Partial Differential Equations (PDEs)}
In physics and engineering, the evolution of complex systems over time and space is governed by Partial Differential Equations (PDEs). For example, the spreading of heat across a lithium-ion battery (Thermal Runaway) is governed by the Heat Equation:
\begin{equation}
\frac{\partial u}{\partial t} = \alpha \nabla^2 u
\end{equation}
Where:
\begin{itemize}
    \item $u(x, t)$ is the temperature at spatial coordinate $x$ and time $t$.
    \item $\alpha$ is the thermal diffusivity of the material.
    \item $\nabla^2$ is the Laplace operator, representing spatial diffusion.
\end{itemize}

Traditionally, finding $u(x, t)$ is impossible to do analytically for complex battery shapes. Engineers rely on \textbf{Finite Element Analysis (FEA)} or Finite Difference Methods (FDM). These methods divide the physical space into a massive grid of tiny discrete points and iteratively solve the math for every single point, moving forward in tiny increments of time (e.g., $\Delta t = 0.001$ seconds). This is computationally grueling. Simulating 10 minutes of battery thermal runaway can take hours or days on a supercomputer.

\subsection{The Failure of Standard Deep Learning}
Deep learning offers a tempting shortcut: train a neural network to look at the initial heat state ($t=0$) and predict the final state ($t=1$). 
However, standard architectures fail fundamentally:
\begin{itemize}
    \item \textbf{Convolutional Neural Networks (CNNs):} CNNs are discretized operators. If you train a ResNet on a $64 \times 64$ sensor grid, it learns fixed $3 \times 3$ pixel filters. If you deploy that model on a factory floor with higher quality $256 \times 256$ sensors, the model completely breaks. It is tied to the resolution. Furthermore, CNNs learn \textit{local} features (edges, corners). Heat diffusion is a \textit{global} phenomenon.
    \item \textbf{Recurrent Neural Networks (LSTMs/Transformers):} Unrolling sequential models across 2D/3D physical space and time results in an explosion of parameters and catastrophic vanishing gradients.
    \item \textbf{Physics-Informed Neural Networks (PINNs):} PINNs bake the PDE directly into the loss function. While mathematically elegant, a standard PINN must be entirely retrained from scratch via gradient descent for every single new initial condition or battery geometry. They are optimization solvers, not generalizable prediction models.
\end{itemize}

\subsection{The Neural Operator Paradigm}
Instead of mapping finite-dimensional vectors to finite-dimensional vectors (like standard NNs), a \textbf{Neural Operator} aims to map \textit{infinite-dimensional function spaces} to \textit{infinite-dimensional function spaces}. 
If we learn the mapping between the functional spaces, we learn the continuous physics itself. Because the learned mapping is continuous, we can query it at any resolution. We can train on low-resolution data and evaluate on high-resolution data with zero retraining (Zero-Shot Super Resolution).

\subsection{Fourier Neural Operators (FNO)}
The FNO is the most successful realization of the Neural Operator paradigm. The core idea is that updating the state of a physical system involves an integral operator (a continuous convolution). Solving convolutions in physical space is mathematically intensive ($\mathcal{O}(N^2)$ complexity). 
However, the FNO bypasses this by moving into the \textbf{Frequency Domain}. By applying the Fast Fourier Transform (FFT), the FNO converts spatial data into a spectrum of wave frequencies. It filters out high-frequency noise, learns physics purely by modifying the dominant low-frequency waves, and transforms back. This reduces the complexity to $\mathcal{O}(N \log N)$ and enables real-time, resolution-invariant PDE solving.

\section{Exhaustive Code Walkthrough}
This section breaks down the PyTorch implementation of the FNO word-by-word.

\subsection{The \texttt{SpectralConv2d} Class}
This is the heart of the FNO, replacing a standard \texttt{nn.Conv2d} layer.

\begin{lstlisting}[language=Python]
class SpectralConv2d(nn.Module):
    def __init__(self, in_channels, out_channels, modes1, modes2):
        super(SpectralConv2d, self).__init__()
        self.in_channels = in_channels
        self.out_channels = out_channels
        self.modes1 = modes1 # Max frequency modes to keep in X dimension
        self.modes2 = modes2 # Max frequency modes to keep in Y dimension
\end{lstlisting}
\textbf{\texttt{modes}}: In a Fourier transform, higher modes represent jagged, high-frequency details. Physics (like heat diffusion) is generally smooth and dominated by low frequencies. By artificially capping the modes we look at, we drastically compress the computation and act as a low-pass filter, forcing the AI to learn the macro-physics rather than memorizing microscopic pixel noise.

\begin{lstlisting}[language=Python]
        # R is a weight tensor that holds Complex Numbers
        self.weights1 = nn.Parameter(
            torch.empty(in_channels, out_channels, self.modes1, self.modes2, dtype=torch.cfloat)
        )
        nn.init.xavier_normal_(self.weights1)
\end{lstlisting}
\textbf{\texttt{torch.cfloat}}: The Fourier domain consists of complex numbers (Real and Imaginary parts). Standard AI uses standard floats. We must explicitly define these weights as complex floats (\texttt{cfloat}) so they can interact with the Fourier coefficients (amplitude and phase).\\
\textbf{\texttt{nn.Parameter}}: Tells PyTorch that these complex weights are learnable via backpropagation.

\begin{lstlisting}[language=Python]
    def forward(self, x):
        batchsize = x.shape[0]
        # 1. Fourier Transform (Space -> Frequency)
        x_ft = torch.fft.rfft2(x)
\end{lstlisting}
\textbf{\texttt{torch.fft.rfft2}}: Real-valued Fast Fourier Transform 2D. We use \texttt{rfft} instead of \texttt{fft} because our input temperatures are strictly real numbers (not complex). The math of \texttt{rfft} exploits Hermitian symmetry to only compute half the spectrum, cutting memory and compute requirements by exactly 50\%.

\begin{lstlisting}[language=Python]
        # 2. Empty tensor for output frequencies
        out_ft = torch.zeros(batchsize, self.out_channels, x.size(-2), 
                             x_ft.size(-1), device=x.device, dtype=torch.cfloat)
        
        # 3. Multiply input frequencies by learnable complex weights
        out_ft[:, :, :self.modes1, :self.modes2] = \
            torch.einsum("bixy,ioxy->boxy", 
                         x_ft[:, :, :self.modes1, :self.modes2], 
                         self.weights1)
\end{lstlisting}
\textbf{Truncation}: We slice the massive tensor and throw away all frequencies higher than \texttt{modes1} and \texttt{modes2}.\\
\textbf{\texttt{torch.einsum("bixy,ioxy->boxy", ...)}}: Einstein Summation Convention. This is a highly optimized way to write matrix multiplications. We take the input tensor \texttt{bixy} and multiply it by the weights \texttt{ioxy}. Sum over the input channels $i$ to produce a final tensor of shape \texttt{boxy} (batch, output\_channels, x, y). This is the core learning step!

\begin{lstlisting}[language=Python]
        # 4. Inverse Fourier Transform (Frequency -> Space)
        x_out = torch.fft.irfft2(out_ft, s=(x.size(-2), x.size(-1)))
        return x_out
\end{lstlisting}
\textbf{\texttt{irfft2}}: Inverse Real Fast Fourier Transform. Now that we have scaled the frequencies using our learned weights, we translate the complex waves back into a physical, real-valued temperature map.

\section{Mathematical Derivations \& Solved Example}

\subsection{The Math Behind the Architecture}
Let $D \subset \mathbb{R}^d$ be a bounded, open spatial domain (the battery). We want to learn an operator $\mathcal{G} : \mathcal{A} \rightarrow \mathcal{U}$ mapping from initial conditions $a(x)$ to solutions $u(x)$.

The iterative update for a hidden representation $v_t$ in a Neural Operator is defined as a continuous integral transform:
\begin{equation}
v_{t+1}(x) = \sigma \left( W v_t(x) + \int_D \kappa(x, y) v_t(y) dy \right)
\end{equation}

Where:
\begin{itemize}
    \item $W$ is the local linear transform (the bypass layer).
    \item $\kappa(x, y)$ is a kernel function (a Green's function).
\end{itemize}

Solving the integral $\int \kappa(x, y) v_t(y) dy$ is intractable for massive 2D grids. We enforce that $\kappa(x, y) = \kappa(x - y)$. This makes the integral a \textbf{Convolution}.

\begin{equation}
\int_D \kappa(x - y) v_t(y) dy = (\kappa * v_t)(x)
\end{equation}

Now, we invoke the \textbf{Convolution Theorem}. The Fourier transform of a convolution of two functions is the pointwise product of their Fourier transforms:
\begin{equation}
\mathcal{F}(\kappa * v_t) = \mathcal{F}(\kappa) \cdot \mathcal{F}(v_t)
\end{equation}

Let $R = \mathcal{F}(\kappa)$. $R$ becomes the set of complex weights our neural network learns. Therefore, the integral operator simplifies entirely to:
\begin{equation}
\mathcal{F}^{-1} \left( R \cdot \mathcal{F}(v_t) \right)
\end{equation}

\subsection{A Fully Solved 1D Mathematical Example}
Let's manually execute the math for a 1D thermal wave on a metal rod.

\textbf{The Setup:}
\begin{itemize}
    \item Space: A rod of length $L = 1$.
    \item Initial Heat State: $v(x) = 3\sin(2\pi x) + 1\sin(4\pi x)$
\end{itemize}
This means the heat is comprised of two waves: a large low-frequency wave (amplitude 3, frequency 1) and a small higher-frequency wave (amplitude 1, frequency 2).

\textbf{Step 1: The Fourier Transform $\mathcal{F}(v)$}\\
The FFT analyzes $v(x)$ and extracts its frequency components (modes).
\begin{itemize}
    \item Mode $k=1$: Amplitude = 3
    \item Mode $k=2$: Amplitude = 1
    \item Mode $k \ge 3$: Amplitude = 0
\end{itemize}
Our frequency tensor $\mathcal{F}(v)$ is conceptually: $[0, 3, 1, 0, 0, 0\dots]$.

\textbf{Step 2: Truncation and Multiplication by $R$}\\
Assume our FNO is configured with \texttt{modes1 = 1}. This means we filter out all frequencies higher than $k=1$. The truncated tensor is $[0, 3]$.

Now, we multiply by the learned AI weight tensor $R$. The AI has been trained on physics data and learned that over 10 minutes, a wave of frequency $k=1$ loses 50\% of its heat.
\begin{itemize}
    \item $R = [0, 0.5]$
    \item Pointwise Multiplication: $[0, 3] \cdot [0, 0.5] = [0, 1.5]$
\end{itemize}

\textbf{Step 3: The Inverse Fourier Transform $\mathcal{F}^{-1}$}\\
We take the manipulated frequency tensor $[0, 1.5]$ and perform the inverse transform back into physical space. The frequency is $k=1$ and the new amplitude is $1.5$.

\textbf{Final Predicted State:}
\begin{equation}
v_{future}(x) = 1.5\sin(2\pi x)
\end{equation}

In three simple steps, the FNO perfectly predicted the future heat state of the rod, completely ignoring calculations at hundreds of microscopic physical points!

\end{document}
